
\documentclass[a4paper,addpoints
	%,answers
	]{exam}
\usepackage{pgf}
\usepackage{underscore}
\usepackage{amssymb}
\usepackage{tabularx}
%\usepackage{fontspec}
%\usepackage{unicode-math}
\usepackage{color}
\usepackage{hyperref}
\usepackage{listings}
\usepackage{amsmath}
\usepackage{amsfonts}
\usepackage{tikz}
\usetikzlibrary{automata, positioning, arrows}

%\setmainfont[Ligatures=TeX]{TeX Gyre Pagella}
%\setsansfont{TeX Gyre Heros}
%\setmonofont{TeX Gyre Cursor}
%\setmathfont{TeX Gyre Pagella Math}



\definecolor{codegreen}{rgb}{0,0.6,0}
\definecolor{codegray}{rgb}{0.5,0.5,0.5}
\definecolor{codepurple}{rgb}{0.58,0,0.82}
\definecolor{backcolour}{rgb}{0.95,0.95,0.92}

\lstdefinestyle{mystyle}{
    backgroundcolor=\color{backcolour},   
    commentstyle=\color{codegreen},
    keywordstyle=\color{magenta},
    numberstyle=\tiny\color{codegray},
    stringstyle=\color{codepurple},
    basicstyle=\ttfamily\footnotesize,
    breakatwhitespace=false,         
    breaklines=true,                 
    captionpos=b,                    
    keepspaces=true,                 
    numbers=left,                    
    numbersep=5pt,                  
    showspaces=false,                
    showstringspaces=false,
    showtabs=false,                  
    tabsize=2
}

\newcommand\given[1][]{\:#1\vert\:}
\DeclareMathOperator{\SPAN}{span}

% SHOW SOLUTION
%\renewcommand{\solutiontitle}{\noindent\textbf{\underline{Answer:}}\par\noindent}
\printanswers
\printanswers
\cfoot{Page \thepage}



\begin{document}
\begin{titlepage}
	\begin{center}

		\textbf{\LARGE Assignment 4} \\

		\textbf{\Large alfar24@ru.is, hlynurh24@ru.is} \\
	\end{center}
	\vspace{10pt}
	\begin{center}
		\includegraphics[height=0.4\textheight]{Reykjavik_University_Logo.png}
	\end{center}
	\vspace{10pt}
	\begin{center}
		\textbf{\Large T-519-STOR: Theory of Computation}
	\end{center}
\end{titlepage}
\medskip
\pagebreak

\section*{Problem 1: the complexity class NP (30 points)}
\begin{questions}
	\question[20] Assume $A \in$ NP and $B \in$ P
	\begin{solutionbox}{\stretch{-1}}
	\end{solutionbox}
\end{questions}

\end{document}
